\chapter{致谢}\label{front:acknowledgements}
\begingroup\small\setlength{\parskip}{2pt}
\begin{center}
  {\color{Teal}\rule{.22\textwidth}{.6pt}}\par
  \vspace{3mm}
  {\itshape 每一份清晰的理解，都建立在前人的思考与分享之上。}
\end{center}
\vspace{3mm}

本书从概念、证明、计算与解题方法四个方面整理概率论与数理统计。
在此感谢教材与讲义作者留下的系统阐述，也感谢数学学习者对严谨表达、直观解释和多种解法的持续追求。
这些积累使知识能够被反复理解、检验与传递。

\section*{参考资料与感谢}
\begin{enumerate}[label=\arabic*.,leftmargin=1.8em,itemsep=4pt]
\item \textbf{茆诗松、程依明、濮晓龙、倪菁《概率论与数理统计教程》}。
感谢其对概率模型、分布、数字特征与统计推断的系统组织，为本书概念主线提供参照。
\item \textbf{李东风《概率统计讲义》}。
感谢其对常见模型、抽样分布、参数估计、假设检验及回归方法的清晰讲解。
\item \textbf{Hao Wu（吴昊）《Introduction to Probability Theory》课程讲义}。
感谢其从可测性、积分、收敛与条件期望出发阐明概率论的理论结构。
\item \textbf{Athanasios Papoulis、S.~Unnikrishna Pillai《概率、随机变量与随机过程》}。
感谢其将概率模型、随机变量变换与应用背景相结合的叙述方式。
\end{enumerate}

感谢\LaTeX、CTeX、TikZ/PGF及相关开源社区对数学表达与书籍排版的支持。
封面配图采用Till Tantau的\href{https://texample.net/title-graphics/}{\textit{Title graphics}}数学分支图，感谢作者与TeXample社区的开放分享。
也感谢每一位认真阅读、独立演算并愿意反馈疑问的读者；对定义、条件、证明和解答的讨论，都是继续修订本书的动力。

\section*{交流与勘误}
欢迎就题目解法、公式推导与排版提出建议。反馈时可注明章节、页码或题号，便于定位与复查。
\begin{tcolorbox}[colframe=Teal!65,colback=Teal!3,boxrule=.45pt,title={联系方式}]
\textbf{作者}\quad 雪落长安无声\par\smallskip
\textbf{知乎}\quad
\href{https://www.zhihu.com/people/black-despair}{https://www.zhihu.com/people/black-despair}\par\smallskip
\textbf{QQ邮箱}\quad
\href{mailto:3547908579@qq.com}{3547908579@qq.com}
\end{tcolorbox}
\vfill
\begin{flushright}\itshape 雪落长安无声\par 2026年10月31日\end{flushright}
\endgroup
